
\subsection{Aveuglement sous clés malveillantes}

\subsubsection{Définition}

\begin{definition}[Blindness]
Un schéma de signature aveugle est dit aveugle sous clés malveillantes si pour tout adversaire PPT
$\mathcal A$,
\[
\operatorname{Adv}^{\mathrm{blind}}_{\mathcal A}(\lambda)
=
\left|
\Pr\!\left[
\begin{array}{l}
(\mathsf{bvk},m_0,m_1)\gets\mathcal A(1^\lambda),\
b\sample\{0,1\},\\
(\rho_{1,i},st_i)\gets\mathsf{BSUser}(\mathsf{bvk},m_i)
\quad(i\in\{0,1\}),\\
(\rho_{2,b},\rho_{2,1-b})\gets
\mathcal A(\rho_{1,b},\rho_{1,1-b}),\\
\sigma_i\gets
\mathsf{BSDerive}(st_i,\rho_{2,i})
\quad(i\in\{0,1\}),\\
\text{if }\exists i,\
\mathsf{BSVerify}(\mathsf{bvk},m_i,\sigma_i)=0:\\
\qquad\text{then }\sigma_0=\sigma_1=\bot,\\
b=\mathcal A(\sigma_0,\sigma_1)
\end{array}
\right]
-\frac12
\right|
= \operatorname{negl}(\lambda).
\]
Autrement dit, aucun adversaire PPT $\mathcal A$ ne peut distinguer si la signature obtenue correspond au message $m_0$ ou au message $m_1$, avec un avantage non négligeable.
\end{definition}

\subsubsection{Preuve: Blindness du schéma de signature}
On veut montrer que, pour deux messages $m_0$ et $m_1$, un signataire malveillant $\mathcal{A}$ ne peut pas distinguer lequel des deux messages a été signé par l'utilisateur. Autrement dit, les vues de $\mathcal{A}$ dans les deux exécutions sont indistinguables :
\[
\operatorname{View}_{\mathcal{A}}(m_0) \approx \operatorname{View}_{\mathcal{A}}(m_1).
\]
Cela signifie que, pour tout adversaire PPT $\mathcal{A}$, l'avantage de distinguer le bit $b$ choisi par l'utilisateur est négligeable :
\[
\left|\Pr[b'=b]-\frac{1}{2}\right|=\operatorname{negl}(\lambda), \text{ où $b'$ est le bit deviné par l'adversaire.}
\]
Pour montrer cette propriété, nous construisons une suite de jeux.

\bigskip

\noindent\textbf{G$_0$:}

Le premier jeu correspond à l'exécution réelle du protocole avec le message $m_0$. L'utilisateur génère les clés OT $(ek_{i,0}, ek_{i,1})$ pour chaque bit du message $m_0$ et les envoie au signataire. Il obtient ensuite une réponse du signataire, effectue le déchiffrement et produit la signature.

\transition{ \textit{On remplace les preuves NIZK par des preuves simulées produites par un simulateur de type Subversion Zero-Knowledge (sub-ZK).}}

\noindent\textbf{G$_1$:}

Ce jeu est identique à G$_0$, à l'exception que les preuves NIZK sont désormais simulées. Par définition de la propriété sub-ZK, même si les paramètres sont choisis (de manière malveillante ou non) par $\mathcal{A}$, les preuves simulées sont indistinguables des preuves réelles. Ainsi,
$\text{G}_0 \approx \text{G}_1$.

\transition{\textit{On applique la propriété sub-ZK à la preuve NIZK de randomisation afin
de supprimer la dépendance à $F(M)$ et $t'$.}}

\noindent\textbf{G$_2$:}

Ce jeu est identique à G$_1$, à l'exception de la preuve NIZK utilisée
pour montrer la validité de la randomisation. Cette preuve est désormais
simulée par le simulateur sub-ZK. La randomisation réelle utilise:
\[
\widehat{\sigma}_1=\sigma_1F(M)^{t'},
\qquad
\widehat{\sigma}_2=\sigma_2g_2^{t'}.
\]
Ici, le simulateur utilise une relation équivalente indépendante
du message :
\[
\widehat{\sigma}_1=\sigma_1F(\mathbf{0})^{t_0},
\qquad
\widehat{\sigma}_2=\sigma_2g_2^{t_0},
\]
Car d'après la propriété de randomisation de Waters démontrée dans
la section \ref{sec:randomisation-waters}, l'aléa final $t+t'$ est
uniformément distribué dans $\mathbb{Z}_p$ et peut être remplacé par un
nouvel aléa $t_0\sample\mathbb{Z}_p$ indépendant.
Le remplacement de $F(M)$ par $F(\mathbf{0})$ est justifié par une 
réduction à l'hypothèse Decisional Diffie--Hellman (DDH),
dont la preuve est présentée dans la section~\ref{sec:DDH-waters}.
La signature finale produite par le simulateur devient alors :
\[
\sigma=(F(\mathbf{0}),\widehat{\sigma}_1,\widehat{\sigma}_2).
\]
Comme la preuve simulée est indistinguable d'une preuve réelle, 
la vue du signataire reste indistinguable entre les deux jeux.
Ainsi $\mathrm{G}_1\approx\mathrm{G}_2$.

\transition{\textit{On simule les composantes $\sigma_{1,i}$ de manière à les rendre indépendantes du message tout en conservant leur produit égal à $\sigma_1$.}}

\noindent\textbf{G$_3$:}

Dans le jeu précédent, la valeur $\sigma_1$ a été simulée de manière
indépendante du message grâce au remplacement de $F(M)$ par $F(\mathbf{0})$.
Il reste toutefois à enlever la dépendance au message dans chacune des
composantes $\sigma_{1,i}$, car la signature vérifie
\[
\sigma_1=\prod_{i=1}^{\ell}\sigma_{1,i}.
\]
Le simulateur choisit donc uniformément
$\sigma_{1,1},\ldots,\sigma_{1,\ell-1}$, puis définit 
\[
\sigma_{1,\ell}
=
\sigma_1
\left(
\prod_{i=1}^{\ell-1}\sigma_{1,i}
\right)^{-1}.
\]
Ainsi,
$
\prod_{i=1}^{\ell}\sigma_{1,i}
=
\sigma_1,
$
et la relation est préservée. Comme $\sigma_1$ est déjà indépendante du message dans G$_2$,
la distribution des composantes simulées
$\sigma_{1,1},\ldots,\sigma_{1,\ell}$ l'est également. Les jeux $\mathrm{G}_2$ et $\mathrm{G}_3$ 
sont donc identiquement distribués, et donc
$\mathrm{G}_2 \approx \mathrm{G}_3$.

\transition{\textit{On utilise la propriété OT hiding pour montrer que la distribution des clés OT ne révèle aucune information sur le message.}}

\noindent\textbf{G$_4$:}

En effet, si $M_i=0$,
$(ek_{i,0},ek_{i,1})=\left(g_1^{y_i},\,g_1^{x_i-y_i}\right)$,
et si $M_i=1$,
$(ek_{i,0},ek_{i,1})=\left(g_1^{x_i-y_i},\,g_1^{y_i}\right).$
Comme $y_i$ est choisi uniformément dans $\mathbb{Z}_p$, la valeur $x_i-y_i$ l'est également. Les deux distributions sont donc identiques et ne révèlent aucune information sur les bits $M_i$. Ainsi,
$\text{G}_3 \approx \text{G}_4$.

\transition{\textit{On remplace le message $m_0$ par le message $m_1$.}}

\noindent\textbf{G$_5$:}

Ce jeu correspond à l'exécution réelle du protocole avec le message $m_1$.\\\\

\noindent\textbf{Conclusion}

Par les transitions effectuées entre les différents jeux, on obtient :
\[
\text{G}_0 \approx \text{G}_1 \approx \text{G}_2 \approx \text{G}_3 \approx \text{G}_4 \approx \text{G}_5.
\]
Ainsi, la vue du signataire malveillant dans l'exécution du protocole avec le message $m_0$ est indistinguable de celle obtenue dans l'exécution avec le message $m_1$ :
\[
\operatorname{View}_{\mathcal{A}}(m_0) \approx \operatorname{View}_{\mathcal{A}}(m_1).
\]
Par conséquent, l'adversaire $\mathcal{A}$ ne peut pas distinguer quel message a été signé par l'utilisateur. Le protocole satisfait donc la propriété de blindness.


\subsubsection{Preuve: randomisation de Waters ($G_1 \rightarrow G_2$)}
\label{sec:randomisation-waters}

La signature obtenue après le déchiffrement est une signature de Waters
de la forme :
\[
\sigma_1 = h_s^a \cdot F(M)^t,
\qquad
\sigma_2 = g_2^t,
\]

où $t\leftarrow\mathbb{Z}_p$ est un aléa choisi lors de la génération de la
signature.
Lors de l'étape de randomisation, l'utilisateur choisit un nouvel aléa
$t'\leftarrow\mathbb{Z}_p$ et calcule :

\[
\begin{aligned}
\widehat{\sigma}_1 &= \sigma_1 F(M)^{t'}
= h_s^a F(M)^t F(M)^{t'}
= h_s^a F(M)^{t+t'},\\
\widehat{\sigma}_2 &= \sigma_2 g_2^{t'}
= g_2^t g_2^{t'}
= g_2^{t+t'}.
\end{aligned}
\]

Or $t'$ est choisi uniformément dans $\mathbb{Z}_p$, donc
$t+t'$ est également uniformément distribuée dans $\mathbb{Z}_p$.
La valeur de l'aléa final $t+t'$ est indépendante de l'aléa
initial $t$ et masque la dépendance du message dans la signature. La
randomisation permet donc de générer une signature dont la distribution ne
révèle aucune information supplémentaire sur le message.
Cette propriété montre que la randomisation d'une signature de Waters
produit une distribution indépendante de l'aléa initial de signature.
Elle permet donc de remplacer, dans la transition entre G$_1$ et G$_2$,
la randomisation réelle par une randomisation simulée utilisant un aléa
indépendant du message.

\subsubsection{Preuve: Réduction DDH ($G_1 \rightarrow G_2$)}
\label{sec:DDH-waters}

Nous montrons que la transition entre les jeux $G_1$ et $G_2$ est
computationnellement indistinguable sous l'hypothèse Decisional
Diffie--Hellman (DDH).
Le simulateur reçoit une instance DDH
$
(g, g^a, g^b, c),
$
où il doit distinguer les deux cas suivants :
\[
c = g^{ab}
\qquad \text{ou} \qquad
c \xleftarrow{\$} G.
\]

À partir de cette instance, le simulateur exécute l'algorithme \textsf{Setup} et construit les paramètres publics ( les bases $u_0, u_1, \ldots$) tel que:
\begin{itemize}
    \item si $c = g^{ab}$, la distribution simulée corresponde au jeu contenant $F(M)$ ;
    \item si $c$ est uniforme dans $G$, la distribution simulée corresponde au jeu contenant $F(\mathbf{0})$.
\end{itemize}

Ainsi, tout adversaire capable de distinguer ces deux jeux permet de résoudre le problème DDH avec un avantage non négligeable.
